\section*{Chapter \ref{ch:qm:convex-optimisation} --- Convex Optimisation}

\begin{solution}{exo:qm:convex-optimisation:1}
Yes: $\sqrt{w^\top\Sigma w} = \lVert L^\top w\rVert$ is a norm composed with a linear map, hence convex. No: the set where the variance is at least a level is the complement of an ellipsoid; the
segment between $w$ and $-w$ passes through $0$, where the variance is zero.
\end{solution}

\begin{solution}{exo:qm:convex-optimisation:2}
First order: $0.0159 \times 0.02 = 0.032\%$, so about $3.180\% + 0.032\% = 3.212\%$. Re-solving at 22\% gives 3.211\%: the multiplier slightly overstates the gain because the price falls as the limit
loosens.
\end{solution}

\begin{solution}{exo:qm:convex-optimisation:3}
Introduce $u, v \ge 0$ with $w - w_0 = u - v$ and require $\mathbf 1^\top(u + v) \le \tau$. At the optimum $u_iv_i = 0$ for every $i$ whenever the limit binds, so $u + v = |w - w_0|$.
\end{solution}

\begin{solution}{exo:qm:convex-optimisation:4}
Stationarity $\Sigma w + \nu\mathbf 1 = 0$ gives $w = -\nu\Sigma^{-1}\mathbf 1$; the budget gives $\nu = -1/(\mathbf 1^\top\Sigma^{-1}\mathbf 1)$, so $w = \Sigma^{-1}\mathbf 1/(\mathbf 1^\top\Sigma^{-1}\mathbf 1)$.
The multiplier's magnitude is the minimum variance, $1/(\mathbf 1^\top\Sigma^{-1}\mathbf 1)$: the rate at which the (half) variance grows with the budget.
\end{solution}

\begin{solution}{exo:qm:convex-optimisation:5}
$\lambda = (1, 1)$ on the two upper bounds and $\nu = -1$ on the equation: $G^\top\lambda + A^\top\nu = (1, 1) - (1, 1) = 0$ and $h^\top\lambda + b^\top\nu = 0.5 - 1 = -0.5 < 0$. In words: $x_1 + x_2 \le 0.5$
contradicts $x_1 + x_2 = 1$.
\end{solution}

\begin{solution}{exo:qm:convex-optimisation:6}
KKT: $\alpha = 2\lambda\Sigma w$, so $w \propto \Sigma^{-1}\alpha$; the binding limit fixes the scale, $w = \sigma\Sigma^{-1}\alpha/\sqrt{\alpha^\top\Sigma^{-1}\alpha}$, and the optimal value is $\sigma\sqrt{\alpha^\top\Sigma^{-1}\alpha}$,
whose derivative in $\sigma$ is $\sqrt{\alpha^\top\Sigma^{-1}\alpha}$, the maximal \omterm{def:qm:estimation:sharpe}{Sharpe ratio}: the \omterm{def:qm:convex-optimisation:strong}{shadow price} of the volatility limit.
\end{solution}

\begin{solution}{exo:qm:convex-optimisation:7}
\texttt{constraint\_costs()}: turnover 0.44\% of utility, gross exposure 0.28\%, position limits 0.19\%, sector neutrality 0.014\%, dollar neutrality 0. The dollar constraint is the sum of the five sector
constraints, so removing it changes nothing; its multiplier in the solver's answer is not unique and says nothing.
\end{solution}

\begin{solution}{exo:qm:convex-optimisation:8}
Dropping constraints in an arbitrary order trades a portfolio that violates a limit chosen by the order of the code, not by anyone's judgement, and hides the conflict. Compute the \omterm{def:qm:convex-optimisation:certificate}{infeasibility
certificate}, report the conflicting rules to the people who own them, and trade nothing new (or a documented fallback) until a limit is changed deliberately.
\end{solution}

\begin{solution}{pb:qm:convex-optimisation:1}
\textbf{1.} Expected return 3.78\%, volatility 3.47\%, utility 3.18\% (annual).
\textbf{2.} Gross exposure (100\%), turnover (20\%), 43 position limits; the neutralities hold as equations.
\textbf{3.} 18 iterations; stationarity and complementarity residuals below $10^{-9}$.
\textbf{4.} 0.0159: one more percent of turnover is worth 1.6 basis points of utility.
\textbf{5.} 0.0185 per unit of gross exposure.
\textbf{6.} From 2.88\% at 5\% to 3.18\% at 20\% and 3.62\% at 80\%.
\textbf{7.} They are the slopes: 0.0210, 0.0184, 0.0159, 0.0119 and 0.0011 at 10\%, 15\%, 20\%, 30\% and 80\%, matching the finite differences of the curve.
\textbf{8.} Turnover 0.44\%, gross 0.28\%, positions 0.19\%, sector neutrality 0.014\%, dollar neutrality 0.
\textbf{9.} It is implied by the sector neutralities (their sum), so it is redundant; multipliers of redundant constraints are not unique.
\textbf{10.} Expected return rises from 3.60\% to 3.92\% between 10\% and 30\% turnover: 1.6 basis points per 1\% of turnover, close to the utility price.
\textbf{11.} The limits are infeasible: phase one ends at $t^* = 0.005$.
\textbf{12.} Weights $1/40$ on each minimum-position row of sector 0 and a multiplier on the sector's neutrality equation; adding the rows gives ``sector 0 at least 20\% long'' against ``sector 0 exactly
0'': the combination $G^\top\lambda + A^\top\nu = 0$ with $h^\top\lambda + b^\top\nu = -0.005 < 0$.
\textbf{13.} As the cone constraint $\lVert L^\top w\rVert \le 0.04$; its \omterm{def:qm:convex-optimisation:strong}{shadow price} is 1.033 ($= \gamma \times 0.04$ with $\gamma = 25.8$), expected return per unit of volatility allowed; the portfolio earns 5.31\%.
\textbf{14.} A matrix with negative eigenvalues makes some portfolios' variance negative, and the optimiser will find and exploit them.
\textbf{15.} In 74 projections it removes all 21 negative eigenvalues at a Frobenius distance of 4.45, changing no entry by more than 0.31.
\textbf{16.} Turnover: it costs the most (0.44\% of utility), and its price is 1.6 basis points per 1\% traded.
\textbf{17.} Only against the transaction costs the limit is there to save: if trading 1\% of capital costs less than 1.6 basis points, the limit is too tight; the limit is a proxy for a cost model that should
be in the objective.
\textbf{18.} Status, KKT residuals, the value of each binding constraint and its multiplier, the ranking by relaxation, and any certificate.
\textbf{19.} Named result: \emph{the constraint that cost the most}: the turnover limit's \omterm{def:qm:convex-optimisation:strong}{shadow price} is 0.0159, 1.6 basis points of utility per 1\% of turnover, and the constraints rank turnover (0.44\%), gross
exposure (0.28\%), position limits (0.19\%), sector neutrality (0.014\%), dollar neutrality (redundant, 0).
\textbf{20.} It prices every limit in the units of the objective, so limits can be argued about with numbers, and it proves infeasibility when there is nothing to trade.
\end{solution}

\begin{solution}{iq:qm:convex-optimisation:1}
Because every local optimum is global, so a solver that satisfies local optimality conditions has the answer, and those conditions (KKT) can be checked; duality then gives bounds, prices
and \omterm{def:qm:convex-optimisation:certificate}{infeasibility certificates}, and polynomial-time algorithms exist.
\iqlookfor{Local equals global, checkable optimality, and duality.}
\end{solution}

\begin{solution}{iq:qm:convex-optimisation:2}
The optimal dual variable of a constraint: the rate at which the optimal objective improves per unit of relaxation. Show the committee the multiplier of the risk or turnover limit in return per
unit, compare it with what the limit protects against (cost, risk), and confirm by re-solving with the limit relaxed.
\iqlookfor{Marginal value in objective units, and confirmation by relaxation.}
\end{solution}

\begin{solution}{iq:qm:convex-optimisation:3}
Feasibility of the primal and of the dual ($\lambda \ge 0$), stationarity of the \omterm{def:qm:convex-optimisation:dual}{Lagrangian}, and $\lambda_ig_i(x) = 0$. In a portfolio: a limit that does not bind has zero price, and a limit with a positive price
is exactly at its bound; positions strictly inside their limits have a marginal utility equal to the prices of the binding constraints they touch.
\iqlookfor{All four conditions and the economic reading of complementarity.}
\end{solution}

\begin{solution}{iq:qm:convex-optimisation:4}
Solve a phase-one (or self-dual) problem and read the Farkas certificate: the constraints with positive weight form a conflicting subset, which can be reported to their owners. Alternatively
minimise the total violation with weights reflecting priorities.
\iqlookfor{Certificates rather than trial and error.}
\end{solution}

\begin{solution}{iq:qm:convex-optimisation:5}
\omterm{def:qm:convex-optimisation:ipm}{Interior-point methods} give high accuracy in a few dozen iterations whatever the size, with exploitable structure (factor covariance) in the linear algebra; first-order methods such as ADMM have
cheap iterations, warm-start well from yesterday's solution and reach trading accuracy quickly. For a daily rebalance with a \omterm{def:qm:covariance-estimation-and-random-matrices:factor}{factor model} either works; ADMM with warm starts is attractive
intraday, interior points when duals must be accurate.
\iqlookfor{Accuracy versus warm starts, and exploiting the factor structure.}
\end{solution}

\begin{solution}{iq:qm:convex-optimisation:6}
Tracking error: $\lVert L^\top(w - b)\rVert_2 \le \mathrm{TE}$ with $\Sigma = LL^\top$, a second-order cone. Impact $|x|^{3/2} \le s$: split $x$ into positive parts and require $x^2 \le rs$ and $r^2 \le x$ with
$r \ge 0$, two rotated cones; sum the $s$ in the objective.
\iqlookfor{The Cholesky trick and the power-cone decomposition.}
\end{solution}
